% =========================================================
% CSC311 Final Project Report Template
% Single column, 12pt, 8 pages max.
% =========================================================

\documentclass[12pt]{article}

% ---------- Packages ----------
\usepackage[margin=1in]{geometry}
\usepackage{parskip}         % blank line between paragraphs, no indent
\usepackage{graphicx}
\usepackage{booktabs}
\usepackage{amsmath, amssymb}
\usepackage{siunitx}
\usepackage[hidelinks]{hyperref}
\usepackage{caption}
\usepackage{enumitem}
\usepackage{xcolor}
\usepackage{mathpazo}

%--------- Headers and footers ------
\usepackage{fancyhdr}
\pagestyle{fancy}
\setlength{\headheight}{14.49998pt}
\fancyhf{}
\fancyhead[L]{CSC311 Project Report}
\fancyhead[R]{Group Members}
\fancyfoot[C]{\thepage}

% ---------- Captions & Lists ----------
\captionsetup{font=small, labelfont=bf}
\setlist[itemize]{noitemsep, topsep=2pt}
\setlist[enumerate]{noitemsep, topsep=2pt}

% ---------- Helpful macros ----------
\newcommand{\bestmodel}{\textit{[Best Model Name]}}
\newcommand{\projacc}{\textbf{[XX.XX\%]}}

% ---------- Safe figure helper ----------
\makeatletter
\newcommand{\maybeincludegraphics}[3]{%
  \begin{figure}[h]
    \centering
    \IfFileExists{#1}{%
      \includegraphics[width=0.7\linewidth]{#1}%
    }{%
      \fbox{%
        \parbox[c][0.28\textheight][c]{0.7\linewidth}{%
          \centering
          \vspace{0.5em}
          \textit{[Placeholder for \texttt{#1}]}\\
          Add your figure here.
          \vspace{0.5em}
        }%
      }%
    }%
    \caption{#2}
    \label{#3}
  \end{figure}%
}
\makeatother

% =========================================================
% Title Page
% =========================================================
\begin{document}

\begin{titlepage}
    \centering
    \vspace*{3cm}

    {\Large {CSC311: Introduction to Machine Learning} \par}
    \vspace{1cm}
    
    {\Large {Project Report} \par}
    \vspace{1cm}
    
    {\large \textit{Optional Creative Project Title} \par}
    \vspace{3cm}
    
    {\large Submitted by \par}
    
%    \vspace{0.5cm}
%    {\large Group Number}
    
    \vspace{0.5cm}
    {\large Student 1, Student 2, Student 3, Student 4 \par}
    
    \vfill
    {\large University of Toronto \\ Winter 2026 \par}
\end{titlepage}
\clearpage
\setcounter{page}{1}

\section{Instructions}

\textbf{Remove this section from your final submission.}

Submit a pdf file called `csc311challenge.pdf' that describes your final model,
as well as the steps that you took to develop this model.

Along with the PDF file,
please also submit any other `.py' or `.ipynb' files that you used while developing
your model. This latter files will not be graded, but are used as evidence that you
completed your own work.
The `.py' or `.ipynb' files do not need to be runnable by the TA, and can rely on
external imports (e.g. sklearn) and data sets that you don't submit.
Only your final `pred.py' file needs to be runnable.

To make use of this template, you should remove the instructions from each of the sections below
and replace them with a written description that addresses the required points.

As a reminder, you may not use genAI to generate any of the written (English) components of this project, including this final report.
Any genAI usage for other parts of your project must be cited and submitted with your final report in a separate file called `cs311genAI.pdf'.

\textbf{This report is due at 10:00 pm on March 30.}


\section{Data Exploration}

\textbf{Worth 2 points. No more than 2 pages.} Describe how you explored your data, and the process through
which you determined your input features. How did you end up representing your data?
What else did you try?  We are looking for:
\begin{itemize}
	\item A thorough exploration of the data, similar to what is presented in labs.
	What are the distributions of the features? How do these features correlate with the target?
	\item If you use figures to answer those questions, an explanation of how you are
	interpreting those figures. Take care that your figures can be interpreted
	meaningfully.
	\item A clear and logical description of how you determined your input features,
	with convincing logical or empirical evidence justifying your choice. 
	Important features are not overlooked (e.g., not removed for ``ease'').
	\item A clear description of the way(s) that you are representing the data in
	your models.
	\item The descriptions should be consistent with the `.py' and/or `.ipynb' files that you
	used while developing your model.
	\item A clear description of how you are splitting your data into various sets.
	You may use k-fold cross validation if you'd like, but if you do, describe how
	you are applying that idea.
\end{itemize}


\section{Model Exploration}

\textbf{Worth 1 point. No more than 1 page.} Describe the model(s) that you evaluated/explored.
We are expecting a thorough exploration of at least 3+ family of models\footnote{A "family" of models are
	a group of similar models. For example, linear models of various types are considered to be a part of
	the same family.}, even if you did't ultimately use those models.
We are looking for:
\begin{itemize}
	\item How you are applied these models. You don't need to reiterate how they work.
	Instead, we're looking for a description of the choices
	you made to apply each model to this task: e.g., what features from the ``Data" section
	did you use for each model? What adjustments (if any) did you need to make?
\end{itemize}


\section{Model Choice and Hyperparameters}

\textbf{Worth 3 points. No more than 3 pages.} How did you determine which model(s)
to use in your final `pred.py', as well as an exploration of hyperparameters.
We are looking for:
\begin{itemize}
	\item A convincing explanation of how you ensured that the evaluation metrics for various
	models are comparable (i.e., are you using a consistent test set?)
	\item A clear description of the evaluation metric(s) used to evaluate your model, 
	as well as justification for its use.
	\item A description of the hyperparameters that you are tuning, hyperparameter value
	combinations that you have tried, and the value of the evaluation metric(s)
	for those hyperparameters.
	We are not looking for an exhaustive search of all possible hyperparameter 
	combinations, but there should be enough evidence to demonstrate that your
	hyperparameter choices are reasonable.
	(Note that just saying "X was the best hyperparameter value based on metric Y out of
	the hyperparameter values we tried" is not enough. Please present the values of the
	evaluation metrics for each hyperparameter value to \textbf{show} that your choice was the best one.)
	\item A clear description of what your final model choice looks like in the
	submitted `pred.py' file.
	\item The descriptions should be consistent with the `.py' and/or `.ipynb' files that you
	used while developing your model.
\end{itemize}


\section{Prediction}

\textbf{Worth 1 points. No more than 1 page.} How well would you expect your model to perform on the
test set? We are looking for:
\begin{itemize}
	\item A point estimate for your performance, \textbf{not a range}.
	\item A reasoned explanation of your expected model performance, with empirical
	evidence supporting your explanation. 
	You are not graded on the closeness of your estimate to the final test accuracy,
	but you are graded on your reasoning.
\end{itemize}


\section{Workload Distribution}

Describe what each person in the group contributed
to the project. A 1-2 sentence description of each person's role is sufficient.
Each student's description must be written by that student in order for them
to receive credit for the project.


\section{References}

Use references whenever appropriate (e.g., when you consult textbooks, research papers, online tutorials, or external code). Cite in the text using \verb|\cite{}|. For example, let's cite the two references provided in the \texttt{.bib} file \cite{bishop2006prml,hastie2009esl}. (Replace these with your own references.)

\bibliographystyle{plain}
\bibliography{references}

\end{document}
